\begin{table}[htbp]
\centering
\caption[Momentos estadísticos y correlación de PC1--PC3 de precipitación - Población Global]
{Asimetría ($\gamma_1$) y curtosis ($\gamma_2$) de la PC1 de precipitación, y correlación de Spearman ($r$) entre PC1, PC2 y PC3 (con signo) y la vorticidad relativa a 850 hPa asociada a la fase (con signo, negativa en el Hemisferio Sur), para la Población Global. El signo de la correlación depende de la convención de signo de cada EOF (Sección~\ref{subsec:analisis_pc}); lo relevante es que el lado del patrón dominante en que cae cada ciclón está asociado a su intensidad. Los valores significativos ($p<0.05$) se muestran en negrita.}
\label{tab:stat_global_tp}
\begin{tabular*}{\textwidth}{@{\extracolsep{\fill}}llrrrrr}
\hline
\textbf{Región} & \textbf{Fase} & \textbf{Asimetría ($\gamma_1$)} & \textbf{Curtosis ($\gamma_2$)} & \textbf{PC1} & \textbf{PC2} & \textbf{PC3} \\
\hline
\textbf{SBR} & Incipiente       & 1.12 & 1.58 & \textbf{-0.52} & \textbf{-0.16} & \textbf{-0.26} \\
             & Intensificación & 0.85 & 0.87 & \textbf{-0.52} & \textbf{-0.13} & \textbf{0.46} \\
             & Madurez         & 1.20 & 1.53 & \textbf{0.17} & \textbf{0.65} & -0.04 \\
             & Decaimiento     & 1.49 & 2.18 & \textbf{0.19} & \textbf{0.15} & \textbf{-0.29} \\
\hline
\textbf{LPB} & Incipiente       & 1.34 & 2.85 & \textbf{-0.49} & \textbf{-0.10} & 0.07 \\
             & Intensificación & 1.03 & 1.44 & \textbf{-0.59} & -0.01 & \textbf{-0.33} \\
             & Madurez         & 1.09 & 1.45 & -0.04 & \textbf{0.61} & \textbf{-0.34} \\
             & Decaimiento     & 2.06 & 6.02 & -0.07 & \textbf{-0.24} & \textbf{-0.13} \\
\hline
\textbf{ARG} & Incipiente       & 2.55 & 9.30 & \textbf{-0.51} & \textbf{-0.18} & -0.04 \\
             & Intensificación & 1.52 & 4.06 & \textbf{-0.66} & \textbf{0.26} & \textbf{-0.22} \\
             & Madurez         & 1.74 & 5.40 & \textbf{-0.53} & \textbf{0.43} & -0.04 \\
             & Decaimiento     & 2.50 & 9.59 & \textbf{-0.25} & 0.00 & \textbf{0.13} \\
\hline
\end{tabular*}
\end{table}
